% 数学ノート集・共通テンプレート v2（積分論2026年度第1部の版面を参考）
% A5判。章内で定義・定理・例・注意等の番号を共通化。
% コンパイル： platex -interaction=nonstopmode math_note_template_v2.tex (2回)
%              dvipdfmx math_note_template_v2.dvi
\documentclass[a5paper,10pt,onecolumn]{jsarticle}
\usepackage[dvipdfmx,top=22mm,bottom=20mm,inner=19mm,outer=18mm,headheight=13pt,headsep=6mm]{geometry}
\usepackage{amsmath,amssymb,amsthm,mathtools}
\usepackage{enumitem,fancyhdr,titlesec}
\usepackage[dvipdfmx,bookmarks=true,colorlinks=false]{hyperref}
\usepackage{pxjahyper}

% ----- 差し替える文書情報 -----
\newcommand{\NoteTitle}{複素微分の幾何学}
\newcommand{\NoteSeries}{数学ノート集}

% ----- 文字・組版 -----
% 参考PDFの質素な明朝＋Computer Modernの印象を引き継ぎながら，
% A5判で窮屈にならないよう本文の行間を少し確保する．
\linespread{1.13}
\setlength{\parindent}{1zw}
\setlength{\parskip}{0pt}
\setlength{\abovedisplayskip}{6pt plus 2pt minus 2pt}
\setlength{\belowdisplayskip}{6pt plus 2pt minus 2pt}
\setlength{\abovedisplayshortskip}{3pt plus 2pt minus 1pt}
\setlength{\belowdisplayshortskip}{4pt plus 2pt minus 1pt}
\setlist[enumerate]{topsep=3pt,itemsep=1pt,parsep=0pt,leftmargin=2.4zw}
\titleformat{\section}{\normalfont\normalsize\bfseries}{\thesection}{0.8zw}{}
\titlespacing*{\section}{0pt}{2.0ex plus .4ex}{.9ex}
\titleformat{\subsection}{\normalfont\normalsize\bfseries}{\thesubsection}{0.7zw}{}
\titlespacing*{\subsection}{0pt}{1.5ex plus .3ex}{.6ex}
\raggedbottom

% ----- 柱とページ番号：ページ上部、罫線なし -----
\fancypagestyle{mathnote}{%
  \fancyhf{}%
  \fancyhead[L]{\footnotesize \NoteSeries\ （\NoteTitle）}%
  \fancyhead[R]{\footnotesize \thepage}%
  \renewcommand{\headrulewidth}{0pt}%
  \renewcommand{\footrulewidth}{0pt}%
}
\pagestyle{mathnote}

% ----- 節内通し番号の定義・定理群 -----
\newtheoremstyle{mathnotetheorem}
 {4.5pt}{4.5pt}{\normalfont}{}%
 {\bfseries}{}{0.65zw}%
 {\thmname{#1}\thmnumber{ #2}\thmnote{（#3）}}
\theoremstyle{mathnotetheorem}
\newtheorem{theorem}{定理}[section]
\newtheorem{definition}[theorem]{定義}
\newtheorem{lemma}[theorem]{補題}
\newtheorem{proposition}[theorem]{命題}
\newtheorem{corollary}[theorem]{系}
\newtheorem{example}[theorem]{例}
\newtheorem{remark}[theorem]{注意}
\renewcommand{\proofname}{証明}
\renewcommand{\qedsymbol}{$\square$}
\makeatletter
\renewenvironment{proof}[1][\proofname]{%
  \par\pushQED{\qed}\normalfont\topsep4\p@\@plus2\p@\relax
  \trivlist\item[\hskip\labelsep\bfseries #1]\ignorespaces
}{\popQED\endtrivlist\@endpefalse}
\makeatother

% ----- 使うノートで適宜拡張する記号 -----
\newcommand{\R}{\mathbb{R}}
\newcommand{\norm}[1]{\lVert #1\rVert}
\newcommand{\abs}[1]{\lvert #1\rvert}

\begin{document}
\section*{読了後にできること}
本稿を読むための前提は，Cauchy--Riemann方程式，Jacobi行列，正則関数．
本稿の具体的な到達課題は，正則写像の微分を回転と拡大縮小に分解する．

\section{複素構造と実線形写像}

本稿では $\mathbb C$ を実二次元線形空間 $\mathbb R^2$ と同一視し，
$z=x+iy$ と書く．
複素数の乗法が与える実線形写像を幾何学的に記述する．

\begin{definition}
$\mathbb R^2$ 上の線形写像
\[
 J:\mathbb R^2\to\mathbb R^2,\qquad J(x,y)=(-y,x)
\]
を標準複素構造という．行列表示は
$J=\begin{pmatrix}0&-1\\1&0\end{pmatrix}$ であり，
$J^2=-I$ を満たす．
\end{definition}

\begin{lemma}[複素線形性の特徴づけ]\label{lem:complex-linear}
実線形写像 $A:\mathbb R^2\to\mathbb R^2$ について，
以下は同値である．
\begin{enumerate}
\item $A$ は $\mathbb C$ 上線形である．
\item $AJ=JA$ が成り立つ．
\item ある $a,b\in\mathbb R$ に対して
\[
 A=\begin{pmatrix}a&-b\\ b&a\end{pmatrix}
\]
と書ける．
\end{enumerate}
このとき $A$ は複素数 $a+ib$ の乗法に等しい．
\end{lemma}
\begin{proof}
実線形な $A$ が複素線形であることは
$A(iz)=iA(z)$ が全ての $z$ に対して成立することと同値であり，
これは $AJ=JA$ に等しい．
$A=\begin{pmatrix}a&c\\b&d\end{pmatrix}$ として
$AJ$ と $JA$ を計算すると
\[
 AJ=\begin{pmatrix}c&-a\\d&-b\end{pmatrix},\qquad
 JA=\begin{pmatrix}-b&-d\\a&c\end{pmatrix}
\]
なので $AJ=JA$ は $c=-b,d=a$ と同値である．
最後の行列は $x+iy$ を
$(a+ib)(x+iy)$ に写す作用に一致する．
\end{proof}

\begin{proposition}\label{prop:rotation}
$a+ib\ne0$ とすると，
複素数 $a+ib$ の乗法を表す実行列は
正の相似拡大と回転の合成である．
具体的には
\[
 \begin{pmatrix}a&-b\\b&a\end{pmatrix}
 =|a+ib|
 \begin{pmatrix}\cos\theta&-\sin\theta\\
                 \sin\theta&\cos\theta\end{pmatrix},
 \qquad a+ib=|a+ib|e^{i\theta}.
\]
その行列式は $a^2+b^2>0$ である．
\end{proposition}
\begin{proof}
極形式の定義から $a=r\cos\theta$，
$b=r\sin\theta$ と表せる．
行列に代入すれば表示式が得られる．
行列式は $a^2+b^2=r^2>0$ である．
\end{proof}

\section{実微分と複素微分}

\begin{definition}
開集合 $U\subset\mathbb C$ 上の写像
$f:U\to\mathbb C$ が $z_0\in U$ で
\emph{複素微分可能}であるとは，
\[
 f'(z_0)=\lim_{h\to0,\ h\ne0}
 \frac{f(z_0+h)-f(z_0)}{h}
\]
が複素数として存在することをいう．
$U$ の全点で複素微分可能な関数を正則関数という．
\end{definition}

\begin{theorem}[複素微分と実Fr\'echet微分]\label{thm:cr}
$f=u+iv:U\to\mathbb C$ とし，
$z_0=x_0+iy_0\in U$ とする．
次の二条件は同値である．
\begin{enumerate}
\item $f$ は $z_0$ で複素微分可能である．
\item $f:\mathbb R^2\to\mathbb R^2$ は $z_0$ で
実Fr\'echet微分可能であり，その微分 $Df(z_0)$ は
複素線形である．
\end{enumerate}
この場合，偏導関数は
\[
 u_x=v_y,\qquad u_y=-v_x
 \tag{CR}
\]
を満たし，
$f'(z_0)=u_x(z_0)+iv_x(z_0)$ である．
\end{theorem}
\begin{proof}
$f$ が複素微分可能なら
\[
 f(z_0+h)=f(z_0)+f'(z_0)h+o(|h|)
\]
が成立する．$h\mapsto f'(z_0)h$ は
複素線形，従って実線形なので，
これは実Fr\'echet微分可能性と
$Df(z_0)h=f'(z_0)h$ を意味する．
よって $Df(z_0)$ は複素線形である．

逆に実Fr\'echet微分可能で
$Df(z_0)$ が複素線形なら，
補題\ref{lem:complex-linear}により
$Df(z_0)h=ch$ となる $c\in\mathbb C$ が存在する．
\[
 \frac{f(z_0+h)-f(z_0)}h
 =c+\frac{o(|h|)}h\longrightarrow c
\]
なので複素微分可能で $f'(z_0)=c$ である．
実ヤコビ行列は
$\begin{pmatrix}u_x&u_y\\v_x&v_y\end{pmatrix}$
である．補題\ref{lem:complex-linear}から
これが複素線形であることは
$u_x=v_y,u_y=-v_x$ と同値である．
また複素乗法係数は $u_x+iv_x$ である．
\end{proof}

\begin{remark}
Cauchy--Riemann方程式がある点で成立するだけでは，
その点で複素微分可能であることは一般には従わない．
定理\ref{thm:cr}の逆方向には
\emph{実Fr\'echet微分可能性}を仮定している．
ただし $u,v$ が $C^1$ 級であれば実Fr\'echet微分可能なので，
Cauchy--Riemann方程式は複素微分可能性の十分条件となる．
\end{remark}

\begin{example}[Cauchy--Riemann方程式だけでは不十分]
\[
 g(z)=
 \begin{cases}
 \dfrac{x^2y}{x^2+y^2}&(z=x+iy\ne0),\\
 0&(z=0)
 \end{cases}
\]
と定める．このとき原点で
$u_x=u_y=v_x=v_y=0$ が成立するが，
$g$ は原点で複素微分可能でない．
\end{example}
\begin{proof}
$g$ は実数値であり，$x$ 軸上と $y$ 軸上で
恒等的に零である．従って原点で両偏導関数は零で，
虚部も零なのでCauchy--Riemann方程式が成立する．
しかし $x=y=t\ne0$ とすると
$g(t+it)=t/2$ だから
\[
 \frac{g(t+it)-g(0)}{t+it}=\frac1{2(1+i)}
\]
となる．一方，実数方向 $h=t$ では差商が零であるため，
複素微分係数は存在しない．
\end{proof}

\section{Wirtinger微分}

\begin{definition}
$C^1$ 級関数 $f=u+iv$ に対して
\[
 \partial_zf=\frac12(\partial_xf-i\partial_yf),
 \qquad
 \partial_{\bar z}f=\frac12(\partial_xf+i\partial_yf)
\]
と定める．これらをWirtinger微分という．
\end{definition}

\begin{proposition}\label{prop:wirtinger}
$f$ が $C^1$ 級であるとする．
このとき実微分は
\[
 df=(\partial_zf)\,dz+(\partial_{\bar z}f)\,d\bar z
\]
と書ける．さらに $f$ が正則であることと
$\partial_{\bar z}f=0$ が成立することは同値である．
\end{proposition}
\begin{proof}
$dz=dx+i\,dy$，$d\bar z=dx-i\,dy$ と書けば
\begin{align*}
 (\partial_zf)\,dz+(\partial_{\bar z}f)\,d\bar z
 &=\tfrac12(f_x-if_y)(dx+i\,dy)\\
 &\quad+\tfrac12(f_x+if_y)(dx-i\,dy)\\
 &=f_xdx+f_ydy=df.
\end{align*}
$\partial_{\bar z}f=0$ は $f_x+if_y=0$ と同値であり，
$f=u+iv$ に分けると
$u_x=v_y$，$u_y=-v_x$ という
Cauchy--Riemann方程式になる．
$f$ は $C^1$ 級なので定理\ref{thm:cr}から
正則性と同値である．
\end{proof}

\begin{example}
$f(z)=z^2$ は $\partial_{\bar z}f=0$，
$f(z)=\bar z$ は $\partial_zf=0$，
$f(z)=|z|^2=z\bar z$ は
$\partial_zf=\bar z$，
$\partial_{\bar z}f=z$ を満たす．
\end{example}
\begin{proof}
$x,y$ に関する偏微分をそれぞれ計算すればよい．
例えば $|z|^2=x^2+y^2$ について
$f_x=2x$，$f_y=2y$ だから
$\partial_zf=x-iy=\bar z$，
$\partial_{\bar z}f=x+iy=z$ である．
\end{proof}

\section{等角性と向き}

\begin{theorem}[正則写像の局所等角性]\label{thm:conformal}
$f:U\to\mathbb C$ を正則関数とし，
$z_0\in U$ で $f'(z_0)\ne0$ とする．
このとき $Df(z_0)$ は非零の相似拡大と回転の合成であり，
任意の非零の接ベクトル $v,w\in\mathbb R^2$ の間の
向き付き角度を保つ．さらに
\[
 \det Df(z_0)=|f'(z_0)|^2>0.
\]
\end{theorem}
\begin{proof}
定理\ref{thm:cr}から
$Df(z_0)$ は $f'(z_0)$ の複素乗法を表す．
命題\ref{prop:rotation}によって，
これが正の相似拡大と回転の合成であることが従う．
相似拡大はベクトルの長さを同じ比で変え，
回転は内積を保存する．
従って二つのベクトルのなす角を保つ．
行列式は $|f'(z_0)|^2$ で正なので，向きも保つ．
\end{proof}

\begin{remark}
$f'(z_0)=0$ の場合には，
微分 $Df(z_0)$ は零写像である．
その点で微分を使って角度を定義することはできないため，
等角性の主張には $f'(z_0)\ne0$ が必要である．
例えば $f(z)=z^2$ は原点で $f'(0)=0$ である．
\end{remark}

\begin{proposition}[反正則写像の向き]\label{prop:antihol}
$f(z)=\overline z$ は実線形の等長写像であるが，
$\det Df=-1$ なので向きを反転する．
\end{proposition}
\begin{proof}
$f(x+iy)=x-iy$ なので
\[
 Df=\begin{pmatrix}1&0\\0&-1\end{pmatrix}.
\]
これは直交行列で行列式が $-1$ である．
従って長さと角度の絶対値は保存するが向きは反転する．
\end{proof}

\section{局所逆写像}

\begin{theorem}[正則逆写像の局所的存在]\label{thm:local-inverse}
$f:U\to\mathbb C$ を $C^1$ 級正則関数，
$f'(z_0)\ne0$ とする．
このとき $z_0$ の開近傍 $U_0$ と
$f(z_0)$ の開近傍 $V_0$ が存在して
$f|_{U_0}:U_0\to V_0$ は全単射となり，
その逆写像は正則である．
さらに $w=f(z)\in V_0$ に対して
\[
 (f^{-1})'(w)=\frac1{f'(z)}
\]
が成り立つ．
\end{theorem}
\begin{proof}
定理\ref{thm:conformal}により
$\det Df(z_0)=|f'(z_0)|^2>0$ だから
$Df(z_0)$ は実線形な可逆写像である．
実逆写像定理を適用すると，
適切な開近傍 $U_0,V_0$ に対して
$f|_{U_0}$ は $C^1$ 級微分同相となり，
その逆写像 $g:V_0\to U_0$ の実微分は
$Dg(w)=(Df(z))^{-1}$ となる．
$Df(z)$ は複素数 $f'(z)$ の乗法なので，
その逆写像は複素数 $1/f'(z)$ の乗法である．
従って $Dg(w)$ は複素線形であり，
定理\ref{thm:cr}から $g$ は複素微分可能である．
また複素導関数は $1/f'(z)$ である．
\end{proof}

\section{調和関数と正則関数}

\begin{definition}
開集合 $U\subset\mathbb R^2$ 上の実数値
$C^2$ 級関数 $u$ が\emph{調和}であるとは，
\[
 \Delta u=u_{xx}+u_{yy}=0
\]
を満たすことをいう．
\end{definition}

\begin{proposition}[正則関数の実部と虚部]\label{prop:harmonic}
$f=u+iv$ が $C^2$ 級正則関数なら，
$u,v$ はいずれも調和関数である．
\end{proposition}
\begin{proof}
Cauchy--Riemann方程式から
$u_x=v_y$，$u_y=-v_x$ である．
両辺をさらに微分すると
\[
 u_{xx}=(v_y)_x=v_{yx},\qquad
 u_{yy}=-(v_x)_y=-v_{xy}.
\]
$C^2$ 級なので混合偏導関数の交換が可能であり，
$\Delta u=0$ を得る．
同様に
$v_x=-u_y$，$v_y=u_x$ から
$\Delta v=-u_{yx}+u_{xy}=0$ が従う．
\end{proof}

\begin{theorem}[星状領域における調和共役]\label{thm:conjugate}
$U\subset\mathbb R^2$ を星状開集合とし，
$u\in C^2(U,\mathbb R)$ が調和関数であるとする．
このとき $v\in C^2(U,\mathbb R)$ が存在して
$f=u+iv$ は正則となる．
さらに $v$ は加法定数を除いて一意である．
\end{theorem}
\begin{proof}
$U$ の星状中心を $a=(a_1,a_2)$ とする．
$P=-u_y$，$Q=u_x$ と置くと
\[
 \partial_xQ-\partial_yP=u_{xx}+u_{yy}=\Delta u=0.
\]
そこで
\[
 v(x)=\int_0^1
 \bigl(P(a+t(x-a))(x_1-a_1)
       +Q(a+t(x-a))(x_2-a_2)\bigr)\,dt
\]
と定める．$u\in C^2$ なので積分と偏微分を交換できる．
$x_1$ に関する微分は
\begin{align*}
 \partial_{x_1}v(x)
 &=\int_0^1\bigl(
   P(a+t(x-a))\\
 &\qquad+t(x_1-a_1)P_x(a+t(x-a))\\
 &\qquad+t(x_2-a_2)Q_x(a+t(x-a))
 \bigr)\,dt.
\end{align*}
$\partial_xQ=\partial_yP$ を用いると，
被積分関数は
$\frac d{dt}[tP(a+t(x-a))]$ に等しい．
従って積分の基本定理から
$v_{x_1}(x)=P(x)=-u_y(x)$ を得る．
同様に $x_2$ に関して微分すると
被積分関数は
$\frac d{dt}[tQ(a+t(x-a))]$ となり，
$v_{x_2}(x)=Q(x)=u_x(x)$ を得る．
従って $v\in C^2(U)$ である．
従って $u+iv$ はCauchy--Riemann方程式を満たし，
定理\ref{thm:cr}から正則である．

$v_1,v_2$ がともに調和共役なら
$(v_1-v_2)_x=(v_1-v_2)_y=0$ である．
星状領域は道連結なので，
任意の二点を結ぶ折れ線または線分に沿って積分すれば
$v_1-v_2$ は定数である．
\end{proof}

\begin{remark}
一般の非単連結領域では，
閉1形式 $\omega$ が大域的に完全でないことがある．
例えば $\mathbb C\setminus\{0\}$ 上の
$u(z)=\log|z|$ は調和だが，
大域的な調和共役は存在しない．
この障害は $\omega$ の閉曲線に沿う積分の非零性として現れる．
\end{remark}


\section{到達課題の完全解答}
\noindent\textbf{複素微分の幾何学的な意味．}
\begin{proposition}
$f(z)=z^2$ を $\mathbb R^2$ の写像と見たとき，
$z=x+iy$ におけるJacobi行列は
\[
Df(x,y)=\begin{pmatrix}2x&-2y\\2y&2x\end{pmatrix}.
\]
$z\ne0$ ではこれは倍率 $2|z|$ の拡大と，
角度 $\arg z$ の回転の合成である．
\end{proposition}
\begin{proof}
$z^2=(x^2-y^2)+i(2xy)$ なので
実部・虚部を偏微分すると表示行列となる．
$z=|z|(\cos\theta+i\sin\theta)$ とすれば
\[
Df=2|z|
\begin{pmatrix}\cos\theta&-\sin\theta\\
\sin\theta&\cos\theta\end{pmatrix}.
\]
よって結論が従う．
\end{proof}
\begin{corollary}
$z\ne0$ では $f$ は局所微分同相であり，
$z=0$ では微分が零写像となる．
\end{corollary}
\begin{proof}
$\det Df=4|z|^2$ なので非零点では逆関数定理を適用できる．
原点では表示行列が零となる．
\end{proof}
\noindent\textbf{到達確認．}
$f(z)=iz$ のJacobi行列を求めよ．
\par\noindent\textbf{解答．}
$i(x+iy)=-y+ix$ だから
$\left(\begin{smallmatrix}0&-1\\1&0\end{smallmatrix}\right)$
であり，$\pi/2$ 回転を表す．
\end{document}